\section{NN-S018: a counting-hard integral gamma filtration on a torsion-free Grothendieck ring}

\paragraph{Statement.}
For the same integral cubic $X_\Phi$ and $v\ge1$ of NN-S017, write
$N=\#\operatorname{SAT}(\Phi)$ and $I=\ker(K_0(X_\Phi)\to K_0(\mathbb Z))$,
where the augmentation is pullback along the specified point. Then
\begin{equation}\label{eq:S018-ring}
 K_0(X_\Phi)=\mathbb Z\oplus I,\qquad I\simeq\mathbb Z^N,\qquad I^2=0.
\end{equation}
The same description holds after base change to any field. In compatible
integral sphere bases, every Adams operation satisfies
\[
                  \psi^m(z)=m^v z\qquad (z\in I, m\ge1).
\]
The complete exterior-power and gamma operations on the augmentation
ideal are
\begin{align}
 \lambda_t(z)&=1+z\sum_{m\ge1}(-1)^{m-1}m^{v-1}t^m,
                                                    \label{eq:S018-lambda}\\
 \gamma_t(z)&=1+z\sum_{j=1}^{v}(-1)^{j-1}(j-1)!
                       \left\{\begin{matrix}v\\j\end{matrix}\right\}t^j.
                                                    \label{eq:S018-gamma}
\end{align}
The braces denote Stirling numbers of the second kind. On an arbitrary
element $a+z$, $a\in\mathbb Z$, the exterior-power series is
$(1+t)^a\lambda_t(z)$, so these formulas specify the entire augmented
$\lambda$-ring.

Define
\[
 d_i(v)=\gcd_{i\le j\le v}\left((j-1)!
                    \left\{\begin{matrix}v\\j\end{matrix}\right\}\right)
           \quad(1\le i\le v),\qquad d_{v+1}(v)=0.
\]
For the integral gamma filtration of SGA~6, Expos\'e~V, Section~3.10,
\begin{equation}\label{eq:S018-filtration}
 F_\gamma^0K_0=K_0,\qquad F_\gamma^iK_0=d_i(v)I\ (1\le i\le v),\qquad
 F_\gamma^{v+1}K_0=0.
\end{equation}
Thus for $1\le i<v$ its associated graded groups are
\[
 \operatorname{gr}_\gamma^iK_0(X_\Phi)
       \simeq\left(\mathbb Z/(d_{i+1}(v)/d_i(v))\right)^N,
 \qquad
 \operatorname{gr}_\gamma^vK_0(X_\Phi)=(v-1)!I\simeq\mathbb Z^N.
\]
In particular, if $v\ge3$ is odd,
\begin{equation}\label{eq:S018-torsion}
             \operatorname{gr}_\gamma^{v-1}K_0(X_\Phi)
                        \simeq (\mathbb Z/2)^N.
\end{equation}
The original verifier can be padded parsimoniously so that $v$ is odd
and at least three. On that family the dimension over $\mathbb F_2$ of
this finite graded group is \#P-complete, its nonvanishing is
NP-complete, and its vanishing is coNP-complete. The ring $K_0$ is nevertheless torsion-free; over the perfect
fields covered by NN-S017, every ordinary integral Chow group
is torsion-free as well.
The torsion is in the associated graded of this filtration.

For completeness the higher groups have the additive description
\[
                         K_n(X_\Phi)\simeq K_n(\mathbb Z)^{\oplus(N+1)}
                         \quad(n\ge0),
\]
and likewise with a field in place of $\mathbb Z$. This last formula
does not assert that the higher groups of the base are torsion-free.

\paragraph{Classical interfaces.}
We use the unstable motivic equivalence NN-S017. Representability and
multiplicative Bott periodicity for algebraic $K$-theory on a regular
Noetherian base are imported in their precise standard forms:
Panin--Pimenov--R\"ondigs, \emph{On Voevodsky's algebraic $K$-theory
spectrum BGL}, April~17, 2007 version, Lemma~1.2.6,
Corollaries~1.2.7--1.2.8, and Lemmas~1.3.2--1.3.3,
Proposition~1.3.5, printed pp.~6--8. The source copy is
\url{https://www.math.uni-bielefeld.de/lag/man/251.pdf}.
The regular-base representability statement is also recorded in
Schlichting--Tripathi, \emph{Geometric models for higher
Grothendieck--Witt groups in $\mathbb A^1$-homotopy theory},
introduction, equation~(1.1). We use ordinary algebraic $K$-theory,
whose statement there has no hypothesis that two be invertible;
the later hermitian theorem has a different scope.

The projective-bundle computation is SGA~6, Expos\'e~VI,
Theorem~1.1. Exterior powers on $K_0$ are SGA~6, Expos\'e~V,
Example~2.2(b) and Lemma~2.2.1; the change of variables defining
$\gamma_t$ is Proposition~3.2(ii), and the filtration is (3.10.1).
We apply these classical operations to the new counting family;
no new general theorem about $\lambda$-rings is asserted.

\paragraph{The ring, Bott generator, and Adams action.}
Every motivic sphere $(\mathbb P^1)^{\wedge v}$ is a simplicial
suspension, since $(\mathbb P^1,\infty)\simeq S_s^1\wedge\mathbb G_m$.
A pointed suspension has null reduced diagonal: its pinch map to the
wedge of two copies has both projections homotopic to the identity,
and the composite of that wedge with the smash product is zero.
The same holds for a wedge of suspensions. Products of two reduced
classes in any represented multiplicative cohomology theory thus
vanish. Applying $K$-theory to NN-S017 proves $I^2=0$.
Bott periodicity identifies the reduced $K_n$ of each sphere with
$K_n$ of the base. The rational point splits off the base summand.
This proves \eqref{eq:S018-ring} and the higher-group formula,
since $K_0(\mathbb Z)=K_0(k)=\mathbb Z$ for a field $k$.

Here is also an integral way to compute the operation on a sphere.
On $(\mathbb P^1)^v$ let $u_j=[\mathcal O(1)_j]-1$.
The projective-bundle theorem gives
\[
 K_0((\mathbb P^1)^v)=
       \mathbb Z[u_1,\ldots,u_v]/(u_1^2,\ldots,u_v^2).
\]
The quotient map from the product to the pointed smash product
pulls the reduced Bott generator back to $u_1\cdots u_v$, up to
a choice of sign. Its pullback on the reduced rank-one group is
injective. Adams operations on a line bundle have
$\psi^m([L])=[L]^m$, and are ring endomorphisms; these are the
Newton identities for the exterior-power operations. Therefore
\[
 \psi^m(u_1\cdots u_v)=
       \prod_{j=1}^v((1+u_j)^m-1)=m^v u_1\cdots u_v.
\]
This proves the displayed Adams formula on each sphere and hence on $I$.

The naturality used here is precisely the unstable operation theorem
of Riou, \emph{Algebraic $K$-theory, $\mathbb A^1$-homotopy and
Riemann--Roch theorems}, arXiv:0907.2710v2, September~9, 2009,
Theorems~1.1.2 and 2.3.1, Proposition~2.3.2, and the opening of
Section~2.4, printed pp.~4 and 10--11. Theorem~2.3.1 gives the
unique special $\lambda$-ring structure on $\mathbb Z\times\mathrm{Gr}$
in the unstable motivic category inducing the usual exterior
powers on every smooth scheme over a regular base. Section~2.4
extends it by representability to every motivic space; the
proposition makes it compatible with regular base change.
Consequently the sphere equivalence and the quotient map above
commute with the usual exterior-power and Adams operations.
This is an unstable integral statement; no periodic stable Adams
map requiring inversion of $m$ is used. The primary copy is
\url{https://www.imo.universite-paris-saclay.fr/~joel.riou/doc/operations.pdf}.

\paragraph{Exact integral formulas and the filtration.}
In the torsion-free ring $K_0\otimes\mathbb Q$, the Newton identity is
\[
 \log\lambda_t(z)=\sum_{m\ge1}\frac{(-1)^{m-1}}m\psi^m(z)t^m.
\]
Since every positive coefficient is in the square-zero ideal,
the logarithm is just $\lambda_t(z)-1$. Substituting the Adams
formula gives \eqref{eq:S018-lambda}. All coefficients obtained
are integral, so torsion-freeness proves the identity already
over $\mathbb Z$.

SGA~6, V, Proposition~3.2(ii), gives
$\gamma_t(z)=\lambda_{t/(1-t)}(z)$. Write its scalar factor as
$\gamma_t(z)=1+P_v(t)z$. The series formula shows
\[
 P_1(t)=t,\qquad P_{v+1}(t)=t(1-t)P_v'(t).
\]
The Stirling recurrence
$\left\{\begin{smallmatrix}v+1\\j\end{smallmatrix}\right\}
=j\left\{\begin{smallmatrix}v\\j\end{smallmatrix}\right\}
+\left\{\begin{smallmatrix}v\\j-1\end{smallmatrix}\right\}$
then proves \eqref{eq:S018-gamma} by induction. This proves the
finite polynomial formula, not merely a formal asymptotic expansion.

By definition $F_\gamma^i$ is generated by products
$\gamma^{j_1}(z_1)\cdots\gamma^{j_s}(z_s)$, where every $z_b\in I$
and $\sum j_b\ge i$. Any product with two positive-degree factors
vanishes because $I^2=0$. The remaining generators are exactly
the scalar multiples
\[
 (j-1)!\left\{\begin{matrix}v\\j\end{matrix}\right\}I,
                         \qquad i\le j\le v.
\]
Their sum is the gcd multiple $d_i(v)I$. Multiplication by an
arbitrary element of $K_0=\mathbb Z\oplus I$ adds no further
generators, again because $I^2=0$. This proves
\eqref{eq:S018-filtration}. Taking the indicated quotients gives
all associated graded groups; $d_i(v)$ divides $d_{i+1}(v)$.

When $v$ is odd and at least three,
$\left\{\begin{smallmatrix}v\\v-1\end{smallmatrix}\right\}
=\binom v2$ and
$\left\{\begin{smallmatrix}v\\v\end{smallmatrix}\right\}=1$,
so
\[
 d_{v-1}(v)=\gcd\left(\frac{v(v-1)!}{2},(v-1)!\right)
                   =\frac{(v-1)!}{2},\qquad d_v(v)=(v-1)!.
\]
This proves \eqref{eq:S018-torsion}. For example at $v=3$,
$\gamma_t(z)=1+(t-3t^2+2t^3)z$,
$F_\gamma^1=F_\gamma^2=I$, $F_\gamma^3=2I$, and
$\operatorname{gr}_\gamma^2=I/2I$.

\paragraph{Counting consequence and geometric meaning.}
Append constant-zero verifier variables, with their Boolean equation
and the equation setting them to zero, until $v\ge3$ is odd.
At most three variables are needed; all old assignments extend
uniquely, all equations retain bounded degree and locality, and
the integral verifier scheme remains the same disjoint union of
$N$ sections in its enlarged affine space. Thus the new cubic
still satisfies every hypothesis of NN-C027 and NN-S017.
Formula~\eqref{eq:S018-torsion} now identifies the indicated
$\mathbb F_2$-dimension with the original exact satisfying count.
This proves all complexity claims, including membership on the
specified syntactic family. E supplies its source counting time,
with the padding contributing only polynomial overhead.

In Grothendieck's language this is an explicit augmented
representation-style $\lambda$-ring of vector bundles, with every
Schur operation determined by its exterior powers and with its
integral gamma filtration computed exactly. Its rational graded
pieces have only gamma degree $v$, corresponding to Tate twist
$v$, while lower
integral graded pieces can have finite torsion of counting-hard
rank. The underlying $K_0$ stays free, as do the ordinary Chow groups
over the perfect fields just specified.
The calculation cautions against replacing the integral gamma
filtration by its rational realization when extracting parity
or multiplicity information. It does not identify these groups
with GCT's canonical multiplicities, prove PH1, or settle P versus
NP or VP versus VNP. Neither A nor M is used.

